\documentclass[conference]{IEEEtran}
\IEEEoverridecommandlockouts

\usepackage{cite}
\usepackage{amsmath,amssymb,amsfonts}
\usepackage{graphicx}
\usepackage{textcomp}
\usepackage{xcolor}
\usepackage{booktabs}
\usepackage{array}
\usepackage{url}

\def\BibTeX{{\rm B\kern-.05em{\sc i\kern-.025em b}\kern-.08em
    T\kern-.1667em\lower.7ex\hbox{E}\kern-.125emX}}

\begin{document}

\title{Passive RF-Silent Drone Detection via Neuromorphic Vision and
Acoustic Sensor Fusion of the Propeller Blade-Pass Frequency}

\author{\IEEEauthorblockN{Giray Yillikci}
\IEEEauthorblockA{\textit{Independent Researcher} \\
gyillikci@gmail.com}}

\maketitle

\begin{abstract}
We investigate passive, RF-silent drone detection by exploiting the fact
that a spinning propeller emits a characteristic blade-pass frequency
(BPF) simultaneously in two physical channels: optical flicker, sensed by
a neuromorphic event camera, and acoustic pressure, sensed by a microphone
array. We present a working prototype that (i)~extracts per-pixel flicker
frequency from a Prophesee EVK4-HD event camera and isolates propeller
regions via connected-component spatial clustering and temporal tracking,
(ii)~extracts the acoustic BPF and a direction-of-arrival (DOA) estimate
from a four-microphone array using multi-channel FFT and GCC-PHAT, and
(iii)~fuses the two streams with a sequential Bayesian updater that
includes harmonic-aware frequency cross-validation. The central finding is
methodological: detection \emph{architecture}, not parameter tuning,
determines success. A grid-FFT baseline failed at range regardless of
tuning, whereas replacing it with per-pixel frequency extraction and
connected-component clustering produced stable, low-false-positive
detection, with first-detection latency reduced from 600--1200\,ms to
approximately 100\,ms after real-time optimization. Acoustic BPF detection
and cross-modal frequency agreement worked reliably; cross-modal
\emph{spatial} (azimuth) calibration did not converge to a usable model
under our indoor test conditions ($R^2 = 0.137$), which we report as an
instructive negative result.
\end{abstract}

\begin{IEEEkeywords}
event camera, neuromorphic vision, drone detection, counter-UAS, sensor
fusion, microphone array, blade-pass frequency, Bayesian inference
\end{IEEEkeywords}

\section{Introduction}
Consumer drones are inexpensive, increasingly autonomous, and---critically
for defense applications---can fly with no active radio-frequency (RF)
link by navigating GPS waypoints. This defeats RF scanning, the dominant
commercial drone-detection technology. Each of the remaining passive and
active modalities has a structural blind spot, summarized in
Table~\ref{tab:modalities}: radar struggles to separate small drones from
birds, frame-based cameras suffer motion blur and limited dynamic range,
and acoustic-only systems exhibit short range and false-positive rates of
10--20\,\% in ambient noise.

\begin{table}[t]
\caption{Structural Weaknesses of Existing Detection Modalities}
\label{tab:modalities}
\centering
\begin{tabular}{@{}p{1.5cm}p{2.6cm}p{3.2cm}@{}}
\toprule
\textbf{Modality} & \textbf{Mechanism} & \textbf{Structural weakness} \\
\midrule
RF scanning & Detect drone radio protocol & Blind to autonomous / RF-silent drones \\
Radar & Reflected RF & Cannot separate small drones from birds; high cost \\
RGB / thermal & Frame imaging (30--60\,fps) & Motion blur, ${\sim}60$\,dB dynamic range, fails in low light / fog \\
Acoustic only & Propeller sound & Short range; 10--20\,\% false-positive rate in noise \\
\bottomrule
\end{tabular}
\end{table}

\subsection{Research Question and Hypothesis}
We ask: can a \emph{passive} sensor combination exploit the
physics-guaranteed propeller blade-pass frequency (BPF) to detect drones
with a false-alarm rate far below acoustic-only systems, without any RF
emission?

Our hypothesis is that because the BPF appears simultaneously in the
optical and acoustic domains, requiring \emph{cross-modal frequency
agreement} should suppress the dominant false-alarm sources---each of
which usually manifests in only one modality---driving the joint
false-positive probability sharply down.

\subsection{Contributions}
This paper makes the following contributions:
\begin{enumerate}
\item A working prototype, to our knowledge the first, that fuses a
      neuromorphic event camera with a microphone array for drone
      detection via the shared BPF observable.
\item An architectural ablation demonstrating that per-pixel frequency
      extraction with connected-component spatial clustering succeeds
      where a grid-FFT baseline fails at range regardless of parameter
      tuning---a transferable lesson for event-based detection pipelines.
\item A harmonic-aware cross-modal frequency validation scheme and a
      gated sequential Bayesian fusion framework, with observed
      accumulate/decay dynamics matching design intent.
\item Real-time engineering techniques (downscaled morphology,
      pre-allocated buffers, region-of-interest statistics) that reduce
      first-detection latency by $6$--$12\times$.
\item An honest account of negative results, including the failure of
      visual-to-acoustic azimuth calibration to converge indoors, and the
      unsuitability of a speech-oriented DOA framework (ODAS) for
      impulsive propeller acoustics.
\end{enumerate}

\section{Related Work}
\textbf{Event-based vision.} Neuromorphic event cameras asynchronously
report per-pixel log-intensity changes with microsecond latency and
$>$120\,dB dynamic range \cite{lichtsteiner2008,gallego2020}. Modern
sensors such as the Sony IMX636 provide HD resolution in a stacked
back-illuminated process \cite{finateu2020}. Their temporal resolution
makes them natural instruments for measuring high-frequency periodic
phenomena such as vibration and rotating machinery---the property we
exploit here.

\textbf{Counter-UAS detection.} Surveys of drone detection cover RF,
radar, electro-optical, and acoustic approaches
\cite{guvenc2018,taha2019}. Acoustic systems detect the propeller's
harmonic signature but suffer short range and high false-alarm rates in
noise; radar and RF systems are active or blind to RF-silent targets. We
are not aware of prior work fusing an event camera with a microphone
array for counter-UAS purposes.

\textbf{Acoustic DOA.} Generalized cross-correlation with phase transform
(GCC-PHAT) \cite{knapp1976} remains the standard lightweight
time-difference-of-arrival estimator; ODAS \cite{grondin2019} provides an
integrated localization/tracking stack tuned primarily for speech.

\textbf{Light-flicker positioning.} A companion line of work uses
luminaire flicker as a location fingerprint \cite{zhang2016litell,
zhu2017ilamp, wang2020vlp}; we discuss a preliminary feasibility study in
Section~\ref{sec:future}.

\section{Physical Principle}
A propeller with $B$ blades rotating at shaft rate $f_s$ (Hz) modulates
both reflected light and radiated sound at the blade-pass frequency
\begin{equation}
f_{\mathrm{BPF}} = B \cdot f_s = \frac{B \cdot \mathrm{RPM}}{60}.
\label{eq:bpf}
\end{equation}
Optically, each blade sweep changes local scene luminance, generating
events in a neuromorphic sensor at $f_{\mathrm{BPF}}$ (and harmonics).
Acoustically, the pressure field carries the same fundamental and its
harmonics. Because the two channels share a single physical cause,
requiring their measured frequencies to agree (up to harmonic ambiguity,
Section~\ref{sec:fusion}) rejects noise sources that are periodic in only
one modality---flickering luminaires, wind noise, speech, or vibrating
structures.

\section{Visual Detection Pipeline}
\label{sec:visual}
The visual pipeline went through nine documented iterations; the crucial
step was an architectural replacement, not tuning.

\subsection{Baseline: Grid-Based FFT (v1) --- Failure at Range}
The first detector divided the $1280\times720$ sensor into a
$16\times16$ grid, counted events per cell per 500\,\textmu s slice, ran a
rolling-window FFT per cell, thresholded on SNR, and clustered adjacent
active cells. This worked only at centimeter range: at distance the
propeller occupies a handful of pixels inside a $160\times90$ cell crowded
with noise pixels, so its spectral peak is diluted below the cell noise
floor.

\subsection{Parameter Tuning --- A Key Negative Result}
To chase range, thresholds were relaxed aggressively (filter length
$7\!\to\!4$, maximum period difference $500\!\to\!1500$\,\textmu s,
minimum pixel count $25\!\to\!5$, grid $16\!\to\!8$ cells, FFT window
$0.5\!\to\!1.0$\,s, minimum SNR $3.0\!\to\!2.0$). The result was 25--43
``propeller regions'' per frame scattered across random
frequencies---pure noise promoted to false positives. The grid FFT has no
notion of spatial coherence: it never asks whether the active pixels in a
cell are spatially or spectrally consistent. \emph{Parameter tuning
cannot repair a fundamentally wrong architecture.}

\subsection{Per-Pixel Frequency + Connected Components (v2)}
The redesigned pipeline is
\begin{equation*}
\begin{split}
\text{per-pixel freq.\ map} &\rightarrow \text{binary mask}
\rightarrow \text{dilation} \\
&\rightarrow \text{connected components} \\
&\rightarrow \text{freq.-consistency filter}
\rightarrow \text{tracking}.
\end{split}
\end{equation*}
Three design choices proved decisive:
\begin{itemize}
\item \textbf{Spatial clustering instead of a fixed grid.} Connected-com\-po\-nent
      labeling adapts to the propeller's true shape; even 2--3 connected
      pixels at a consistent frequency are meaningful, while the same
      pixels scattered across a grid cell are not.
\item \textbf{Frequency-consistency gate.} Within each cluster the
      coefficient of variation of per-pixel frequency must satisfy
      $\mathrm{CV} \le 0.3$.
\item \textbf{Temporal tracking.} Greedy association on spatial
      proximity ($<$100\,px) and frequency similarity ($<$30\,\%);
      confirmation only after a minimum number of consecutive hits;
      pruning after 8 idle frames; exponential smoothing
      ($\alpha = 0.3$) of position and frequency.
\end{itemize}
On a 3-blade bench propeller this yielded a single stable track at
${\sim}99.7$\,Hz ($\approx$2{,}990 blade passes per minute equivalent),
locked position, ${\sim}100$ active pixels, and $270{+}$ consecutive hits
over the recording, with transient single-pixel false positives pruned
within 1--2\,s.

\subsection{Ablation: Minimum Cluster Size}
Setting the minimum cluster size to a single pixel produced 13--17
``confirmed propellers'' per frame at random frequencies, confirming that
\emph{spatial coherence ($\ge$2 connected pixels) is the primary
false-positive filter}: a single pixel carries no spatial information and
is indistinguishable from sensor noise. The final configuration (minimum
2 pixels, 5 confirmation hits) yielded clean detections with zero false
positives and multi-pixel clusters (35--269\,px) tracked over $80{+}$ hits
in a 32\,s recording.

\subsection{Real-Time Optimization (v3)}
The synchronous frequency-map callback initially ran full-resolution
morphology, connected-component analysis, and a frame copy at 25\,Hz,
costing ${\sim}315$\,ms of CPU per wall-clock second and causing
monotonically growing display lag. Four fixes eliminated the drift:
(1)~$4\times$ downscaling ($1280\times720 \to 320\times180$) before
morphology; (2)~reducing the update rate to 10\,Hz; (3)~in-place frame
drawing; and (4)~early exit on near-empty masks, reducing callback cost
to ${\sim}20$\,ms/s. A subsequent production pass added pre-allocated
buffers (eliminating garbage-collection jitter), region-of-interest-only
frequency statistics, velocity-predictive tracking, and a Bayesian
confidence score replacing the hard hit-count threshold: each observation
contributes likelihood
\begin{equation}
p = 0.65 + q\,(0.85 - 0.65),
\end{equation}
where $q \in [0,1]$ combines pixel-count and frequency-consistency
sub-scores, accumulated as
\begin{equation}
P = 1 - (1 - p)^n
\end{equation}
over $n$ observations and decayed by $0.8\times$ per idle frame.
Table~\ref{tab:latency} summarizes the measured effect.

\begin{table}[t]
\caption{Visual Detector Performance, v2 vs.\ v3}
\label{tab:latency}
\centering
\begin{tabular}{@{}lccc@{}}
\toprule
\textbf{Metric} & \textbf{v2} & \textbf{v3} & \textbf{Gain} \\
\midrule
First detect (strong signal) & 600--1200\,ms & ${\sim}100$\,ms & $6$--$12\times$ \\
First detect (weak signal) & 600--1200\,ms & ${\sim}200$\,ms & $3$--$6\times$ \\
Steady-state update & 100--140\,ms & 50--90\,ms & ${\sim}2\times$ \\
Per-frame analysis & ${\sim}1.5$\,ms & ${\sim}0.8$\,ms & ${\sim}2\times$ \\
GC jitter & $\pm 5$\,ms & $<\pm 1$\,ms & ${\sim}5\times$ \\
\bottomrule
\end{tabular}
\end{table}

\section{Acoustic Detection Pipeline}
\label{sec:acoustic}

\subsection{BPF Extraction}
A rolling 0.2\,s window (3200 samples at 16\,kHz) per channel is
Hann-windowed and transformed; magnitude spectra are incoherently averaged
across the four channels (${\sim}6$\,dB SNR gain). The peak is picked in
the target band, its SNR computed against the median noise floor, and
detections below a minimum SNR of 3.0 rejected. This proved reliable on
sustained propeller tones; source calibration on the bench propeller
measured a dominant tone at 162.7\,Hz with SNR $\approx 24$\,dB.

\subsection{DOA via GCC-PHAT}
For each opposing microphone pair (64\,mm baseline) we compute the
phase-transform cross-spectrum
\begin{equation}
R_{12}(\tau) = \int
\frac{X_1(\omega) X_2^{*}(\omega)}{\left| X_1(\omega) X_2^{*}(\omega)
\right|} \, W(\omega)\, e^{j\omega\tau} \, d\omega,
\label{eq:gccphat}
\end{equation}
with a frequency weighting $W(\omega)$ emphasizing the target band,
$4\times$ zero-padded inverse transform for sub-sample delay resolution,
parabolic peak interpolation, and a peak-to-second-peak confidence-ratio
gate ($\ge 1.5$). Pair estimates are fused by averaging when both are
confident, followed by a 7-sample median filter; the sample spread gates
unstable bearings. This yields usable azimuth on steady tones; the planar
32\,mm-radius array cannot resolve elevation and provides
${\sim}\pm 10^{\circ}$ azimuth accuracy at range.

\subsection{ODAS --- A Negative Result}
We integrated the ODAS framework \cite{grondin2019} via a custom socket
relay. It produced robust DOA on speech (sustained utterances of 500\,ms
or more) but was unreliable on propeller acoustics, which are
comparatively impulsive and broadband. This motivated the custom GCC-PHAT
implementation above. The array's on-board DSP additionally exposes a
firmware DOA with speech-tuned voice-activity detection; it was likewise
unreliable on steady tones and was not fused.

\section{Sensor Fusion}
\label{sec:fusion}

\subsection{Harmonic-Aware Frequency Cross-Validation}
One modality frequently reports the blade-pass tone while the other
reports the shaft fundamental or a harmonic. We therefore search harmonic
pairs $(n_v, n_a)$, $1 \le n_v, n_a \le 4$, and declare a match if
\begin{equation}
\frac{\left| f_v / n_v - f_a / n_a \right|}{f_a / n_a} < 0.05,
\label{eq:harmonic}
\end{equation}
returning confidence $1 / (n_v n_a)$ so that fundamental-to-fundamental
agreement outweighs high-order coincidences.

\subsection{Sequential Bayesian Updating}
Starting from prior $P(D) = 0.001$, each modality contributes a
likelihood-ratio update with sigmoid likelihoods:
\begin{align}
P(V \mid D) &= \sigma(\mathrm{SNR}_v - 5), &
P(V \mid \neg D) &= 0.01\, e^{-0.5\,\mathrm{SNR}_v}, \\
P(A \mid D) &= \sigma(\mathrm{SNR}_a - 4), &
P(A \mid \neg D) &= 0.05,
\end{align}
with the acoustic false-alarm likelihood reduced to 0.005 when the noise
floor is elevated. A cross-modal frequency match contributes
$P(F \mid D) = 0.99 \cdot \mathrm{conf}$ against
$P(F \mid \neg D) = 0.001$. When no evidence arrives, the posterior
decays toward the prior by $0.995\times$ per frame. A drone is declared
at $P > 0.8$.

\subsection{Fusion Gating}
To prevent overconfident fusion from coincidental evidence, several gates
proved necessary in practice: a visual quality gate (minimum pixel count
and track confidence $\ge 0.98$); an acoustic stability gate (persistence
and frequency jitter $\le 25$\,Hz); a temporal association gate
($|t_v - t_a| \le {\sim}200$\,ms); an angle-residual gate (visual azimuth
vs.\ acoustic DOA within $35^{\circ}$); and stricter audio-only gates
(SNR $\ge 6$\,dB \emph{and} bounded DOA spread) before an audio-only
update may raise $P(D)$. In a live run with no propeller in view, the
detector repeatedly latched $P(D) \to {\sim}1.0$ on strong 140--200\,Hz
tones through the audio-only path and decayed back when the source
stopped---the accumulate/decay dynamics behaved as designed, and the run
illustrates that single-modality dominance is a real failure mode that
the gates must police.

\section{Implementation}
\textbf{Visual.} Prophesee EVK4-HD (Sony IMX636, $1280\times720$,
${\sim}1$\,\textmu s per-pixel temporal resolution, $>$120\,dB dynamic
range) \cite{finateu2020}, with a $220^{\circ}$ fisheye lens for
hemispheric coverage and a rectilinear C-mount lens
($63.8^{\circ}$ HFOV) for azimuth experiments; Metavision SDK frequency
map algorithms.
\textbf{Acoustic.} ReSpeaker 4-mic circular USB array (32\,mm radius,
64\,mm opposing-pair baseline), 16\,kHz / 16-bit capture; on-board XMOS
XVF-3000 DSP and a 12-LED ring used as a physical bearing indicator.
Pixel-to-azimuth mapping uses a pinhole model,
$f_{\mathrm{px}} = (W/2)/\tan(\mathrm{HFOV}/2)$,
$\mathrm{az} = \mathrm{atan2}(x - W/2,\, f_{\mathrm{px}})$.
\textbf{Software.} Python~3.9 (required by the SDK binaries), NumPy,
OpenCV, PyAudio. The visual Nyquist constraint is
$f_{\max} \le 10^{6} / (2\,\Delta t)$ for accumulation period
$\Delta t$ in \textmu s (1\,kHz at the default
$\Delta t = 500$\,\textmu s). Session traces are logged as JSONL with
per-interval posterior, evidence, association lag, and gate states.

\section{Experiments and Results}
All experiments used a bench propeller as a controlled rotating source
and a synthetic on-screen vibration generator (15--120\,Hz) for
end-to-end validation of the visual frequency pipeline. Recordings were
indoor, single-source, and ${\sim}30$\,s long.

\subsection{Summary of Outcomes}
Table~\ref{tab:outcomes} consolidates what worked and what failed across
the study.

\begin{table}[t]
\caption{Consolidated Outcomes}
\label{tab:outcomes}
\centering
\begin{tabular}{@{}p{5.6cm}c@{}}
\toprule
\textbf{Method / component} & \textbf{Outcome} \\
\midrule
Per-pixel flicker-frequency extraction & \checkmark \\
Connected-component spatial clustering & \checkmark \\
Temporal tracking + Bayesian confidence & \checkmark \\
Real-time latency engineering & \checkmark \\
Acoustic BPF via multi-channel FFT & \checkmark \\
GCC-PHAT azimuth on sustained tones & \checkmark \\
Harmonic-aware frequency cross-validation & \checkmark \\
Gated sequential Bayesian fusion & \checkmark \\
\midrule
Grid-FFT architecture (v1) & $\times$ \\
Parameter tuning to extend range & $\times$ \\
Single-pixel minimum cluster size & $\times$ \\
ODAS for propeller acoustics & $\times$ \\
Visual$\rightarrow$acoustic azimuth calibration & $\times$ \\
Elevation DOA (planar array) & $\times$ \\
\bottomrule
\end{tabular}
\end{table}

\subsection{Azimuth Calibration --- A Negative Result}
\label{sec:azimuth}
To validate the angle-residual gate we attempted to fit a linear map
$\mathrm{DOA}_a = g \cdot \mathrm{az}_v + b$ by sweeping a co-located
source laterally, binning samples by visual azimuth with robust per-bin
medians, and fitting with iterative MAD outlier rejection and an $R^2$
solidity check (threshold 0.6). The best fit (rectilinear lens,
$63.8^{\circ}$ HFOV) achieved $g = -0.011$, $b = 0.37^{\circ}$,
$R^2 = 0.137$, RMSE $0.70^{\circ}$---not solid. The acoustic DOA spanned
only ${\sim}\pm 1.2^{\circ}$ while the visual sweep spanned
${\sim}\pm 25^{\circ}$, consistent with indoor reverberation, an
insufficiently wide or slow sweep, or low GCC-PHAT confidence during the
session. Cross-modal \emph{spatial} fusion therefore remains unvalidated;
cross-modal \emph{frequency} fusion is validated.

\section{Discussion and Lessons Learned}
\textbf{Architecture beats tuning.} The single most consequential change
in the project was architectural: replacing grid-FFT event counting with
per-pixel frequency extraction and connected-component clustering. No
amount of threshold relaxation rescued the grid detector; it merely
promoted the noise floor to false positives.

\textbf{Spatial coherence is the false-positive filter.} The
single-pixel ablation demonstrates that requiring as few as two connected
pixels at a consistent frequency removes essentially all sensor-noise
detections. Persistence (temporal tracking) removes roughly 95\,\% of the
remaining transients.

\textbf{Modality dominance is a real failure mode.} The audio-only run
in Section~\ref{sec:fusion} shows that a Bayesian fuser will happily
saturate on one strong modality; explicit quality, association, and angle
gates---not the update rule itself---are what make fusion trustworthy.

\textbf{Tool-domain mismatch.} ODAS and the array's firmware DOA, both
tuned for speech, failed on propeller acoustics---a reminder that
off-the-shelf audition stacks encode source assumptions that must be
checked before integration.

\section{Limitations and Future Work}
\label{sec:future}
Table~\ref{tab:limits} lists the principal threats to validity. All
detection results are from indoor, short, single-source lab recordings;
no outdoor, multi-drone, or GPS-ground-truth validation was performed,
so range figures are projections rather than measurements, and the
sub-0.1\,\% false-positive-rate target is design-derived and must be
established empirically on a labeled dataset. The 32\,mm planar array
cannot resolve elevation, harmonic ambiguity (a 2-blade propeller at
100\,Hz BPF is indistinguishable from a 4-blade at 50\,Hz shaft rate)
requires external metadata, and there is no hardware inter-sensor clock
(${\sim}100$\,ms association uncertainty).

Future work includes outdoor validation with GPS ground truth, a
non-planar or larger-baseline array for elevation, hardware time
synchronization, multi-target tracking, and repeating the azimuth
calibration outdoors with a wider, slower sweep. A companion feasibility
study on using ceiling-light flicker frequencies as location
fingerprints---inspired by LiTell \cite{zhang2016litell},
iLAMP \cite{zhu2017ilamp}, and event-camera visible-light positioning
\cite{wang2020vlp}---completed a literature review and utility
implementation with a positive feasibility verdict, but end-to-end
validation remains open.

\begin{table}[t]
\caption{Limitations and Threats to Validity}
\label{tab:limits}
\centering
\begin{tabular}{@{}p{2.0cm}p{5.2cm}@{}}
\toprule
\textbf{Area} & \textbf{Limitation} \\
\midrule
Dataset & Indoor, single source, ${\sim}30$\,s recordings \\
Range & No outdoor / GPS ground-truth tests \\
Spatial fusion & Azimuth calibration weak ($R^2 = 0.137$) \\
DOA & Planar 32\,mm array; ${\sim}\pm 10^{\circ}$, no elevation \\
Multi-target & One source per cluster assumed \\
Synchronization & No hardware clock; ${\sim}100$\,ms uncertainty \\
Harmonics & Blade-count / shaft-rate ambiguity \\
FPR claim & Design-derived, not yet measured \\
\bottomrule
\end{tabular}
\end{table}

\section{Conclusion}
We demonstrated a passive, RF-silent drone-detection prototype that fuses
neuromorphic vision and acoustic sensing through the shared blade-pass
frequency observable. The validated contributions are the architectural
comparison (per-pixel frequency plus connected components decisively
outperforms grid FFT), the real-time engineering that brings
first-detection latency to ${\sim}100$\,ms, harmonic-aware cross-modal
frequency validation, and gated sequential Bayesian fusion whose
accumulate/decay dynamics behave as designed. We report the failure of
indoor visual-to-acoustic azimuth calibration as a genuine, instructive
negative result rather than an omission. Range and false-positive-rate
figures remain design targets pending outdoor validation on labeled data.

\begin{thebibliography}{00}
\bibitem{lichtsteiner2008} P. Lichtsteiner, C. Posch, and T. Delbruck,
``A $128\times128$ 120 dB 15 $\mu$s latency asynchronous temporal
contrast vision sensor,'' \emph{IEEE J. Solid-State Circuits}, vol. 43,
no. 2, pp. 566--576, 2008.
\bibitem{gallego2020} G. Gallego \emph{et al.}, ``Event-based vision: A
survey,'' \emph{IEEE Trans. Pattern Anal. Mach. Intell.}, vol. 44, no. 1,
pp. 154--180, 2022.
\bibitem{finateu2020} T. Finateu \emph{et al.}, ``A $1280\times720$
back-illuminated stacked temporal contrast event-based vision sensor with
4.86 $\mu$m pixels, 1.066 GEPS readout, programmable event-rate
controller and compressive data-formatting pipeline,'' in \emph{IEEE Int.
Solid-State Circuits Conf. (ISSCC)}, 2020, pp. 112--114.
\bibitem{guvenc2018} I. Guvenc, F. Koohifar, S. Singh, M. L. Sichitiu,
and D. Matolak, ``Detection, tracking, and interdiction for amateur
drones,'' \emph{IEEE Commun. Mag.}, vol. 56, no. 4, pp. 75--81, 2018.
\bibitem{taha2019} B. Taha and A. Shoufan, ``Machine learning-based drone
detection and classification: State-of-the-art in research,'' \emph{IEEE
Access}, vol. 7, pp. 138669--138682, 2019.
\bibitem{knapp1976} C. Knapp and G. Carter, ``The generalized correlation
method for estimation of time delay,'' \emph{IEEE Trans. Acoust., Speech,
Signal Process.}, vol. 24, no. 4, pp. 320--327, 1976.
\bibitem{grondin2019} F. Grondin and F. Michaud, ``Lightweight and
optimized sound source localization and tracking methods for open and
closed microphone array configurations,'' \emph{Robot. Auton. Syst.},
vol. 113, pp. 63--80, 2019.
\bibitem{zhang2016litell} C. Zhang and X. Zhang, ``LiTell: Robust indoor
localization using unmodified light fixtures,'' in \emph{Proc. ACM
MobiCom}, 2016, pp. 230--242.
\bibitem{zhu2017ilamp} S. Zhu and X. Zhang, ``Enabling high-precision
visible light localization in today's buildings,'' in \emph{Proc. ACM
MobiSys}, 2017, pp. 96--108.
\bibitem{wang2020vlp} G. Wang \emph{et al.}, ``Visible light positioning
using an event camera,'' \emph{IEEE Sensors}, 2020.
\end{thebibliography}

\end{document}
